 Die in dieser Arbeit beschriebenen Tätigkeiten wurden im Rahmen der dritten Praxisphase bei \textbf{ASEA Brown Boveri AG (ABB)} durchgeführt. Der Praxiszeitraum erstreckte sich über die vereinbarte Dauer der dritten Praxisphase. Die Tätigkeiten fanden an den Standorten Frankfurt am Main und Mannheim statt und wurden in der Abteilung \textbf{Process Automation Energy Industries (PAEN)}, konkret im \textbf{Service Zentralbereich Cyber-Security (SZC)}, unter der fachlichen Anleitung der betrieblichen Betreuer Philip Dosch und Jake Tran durchgeführt.
\newline
\newline
Gegenstand dieser Praxisarbeit ist die Analyse moderner \textbf{Command-and-Control-Infrastrukturen (C2)} im Kontext von Red-Team-Operationen. Der Schwerpunkt liegt auf der technischen Betrachtung aktueller C2-Frameworks, wobei insbesondere das Framework \textbf{Sliver} detailliert analysiert wird. Ergänzend erfolgt eine vergleichende Einordnung weiterer verbreiteter Lösungen, darunter unter anderem \textbf{Metasploit} und \textbf{Cobalt Strike}, um Gemeinsamkeiten, Unterschiede sowie jeweilige Einsatzszenarien herauszuarbeiten.
\newline
\newline
Im Rahmen der Arbeit werden zentrale Aspekte moderner C2-Frameworks untersucht, darunter Architekturkonzepte, Kommunikationsmechanismen, Payload-Generierung sowie operative Gesichtspunkte im Hinblick auf Stealth und OPSEC. Ziel ist es, ein fundiertes Verständnis für die Funktionsweise und den praktischen Einsatz dieser Infrastrukturen in realistischen Red-Team-Szenarien zu vermitteln.
\newline
\newline
Die Analyse zeigt zugleich, dass trotz fortschreitender Entwicklung moderner C2-Frameworks weiterhin erkennbare Schwachstellen in Bezug auf Erkennbarkeit, Signaturbildung und Verhaltensanalyse bestehen. Diese Erkenntnisse verdeutlichen die Notwendigkeit zusätzlicher Maßnahmen zur Tarnung von Payloads und Kommunikationsmustern, insbesondere im Umfeld zunehmend leistungsfähiger Endpoint-Detection-and-Response-Systeme.
\newline
\newline
Aufbauend auf den Ergebnissen dieser Praxisphase ergibt sich eine inhaltliche Überleitung zur vierten Praxisphase, in der der Fokus auf der vertieften Auseinandersetzung mit \textbf{Obfuscation-Techniken} liegen wird. Geplant ist dabei die Konzeption und Entwicklung einer eigenen \textbf{Obfuscation-Toolbox}, welche gezielt auf die im Rahmen dieser Arbeit identifizierten Schwachstellen moderner C2-Infrastrukturen eingeht und deren Einsatz in Red-Team-Operationen weiter optimiert.
\newline
\newline
Abschließend möchte ich mich bei meinen Praxisbetreuern, meinem Vorgesetzten sowie bei allen Kolleginnen und Kollegen der Abteilung bedanken, die mich während der dritten Praxisphase fachlich unterstützt haben und jederzeit für einen offenen Austausch zur Verfügung standen.